\section{Lineare und zeitinvariante Systeme}

\subsection{Lineare verschiebungsinvariante Systeme}

Ein LTI/LSI System wird durch seine \textit{Impulsantwort} charakterisiert:

\begin{center}
\begin{tikzpicture}[>=latex, node distance=1.2cm]
  \node (d1) {$\delta(t)$};
  \node [draw, rectangle, right=1cm of d1] (S1) {$\mathcal{S}_{\mathrm{LTI}}[\bullet]$};
  \node [right=1cm of S1] (h1) {$h(t)$};
  \draw[->] (d1) -- (S1);
  \draw[->] (S1) -- (h1);

  \node [below=0.8cm of d1] (d2) {$\delta[k]$};
  \node [draw, rectangle, right=1cm of d2] (S2) {$\mathcal{S}_{\mathrm{LSI}}[\bullet]$};
  \node [right=1cm of S2] (h2) {$h[k]$};
  \draw[->] (d2) -- (S2);
  \draw[->] (S2) -- (h2);
\end{tikzpicture}
\end{center}

Voraussetzung hierfür ist, dass das System vor der Anregung mit dem Impuls in Ruhe ist, d.\,h.\ dass alle inneren Zustände Null sind.

\subsection{Das Faltungsintegral (engl. Convolution)}

Mit der Impulsantwort $h(t) = \mathcal{S}[\delta(t)]$ eines LTI-Systems $\mathcal{S}$ und mit

$$
x(t) = \int_{-\infty}^{\infty} x(\xi)\,\delta(\xi - t)\,\mathrm{d}\xi = \int_{-\infty}^{\infty} x(\xi)\,\delta(t - \xi)\,\mathrm{d}\xi
$$

erhalten wir für ein beliebiges Eingangssignal $x(t)$

\begin{align*}
y(t) = \mathcal{S}\left[x(t)\right] &= \mathcal{S}\left[\int_{-\infty}^{\infty} x(\xi)\,\delta(t-\xi)\,\mathrm{d}\xi\right] \\
&= \int_{-\infty}^{\infty} x(\xi)\,\mathcal{S}\left[\delta(t-\xi)\right]\mathrm{d}\xi \\
&= \int_{-\infty}^{\infty} x(\xi)\,h(t-\xi)\,\mathrm{d}\xi = \int_{-\infty}^{\infty} h(\xi)\,x(t-\xi)\,\mathrm{d}\xi \\
&= x(t) * h(t)
\end{align*}

\begin{beispiel}[RC-Netz: Berechnung des Ausgangssignals]
Die Eingangs-Ausgangsbeziehung des RC-Netzes ist wie folgt gegeben:

$$
y(t) = z(t_0)\,e^{-\frac{t-t_0}{\tau}} + \frac{1}{\tau}\int_{t_0}^{t} x(\xi)\,e^{-\frac{t-\xi}{\tau}}\,\mathrm{d}\xi
$$

Für $z(t_0) = 0$ und $t_0 = -\infty$ erhält man ein Faltungsintegral:

\begin{align*}
y(t) &= \frac{1}{\tau}\int_{-\infty}^{t} x(\xi)\,e^{-\frac{t-\xi}{\tau}}\,\mathrm{d}\xi \\
     &= \int_{-\infty}^{\infty} x(\xi)\,\frac{1}{\tau}\,e^{-\frac{t-\xi}{\tau}}\,\epsilon(t-\xi)\,\mathrm{d}\xi \\
     &= \int_{-\infty}^{\infty} x(\xi)\,h(t-\xi)\,\mathrm{d}\xi
\end{align*}

$\Rightarrow$ Kausales System mit der Impulsantwort $h(t) = \epsilon(t)\,\dfrac{1}{\tau}\,e^{-\frac{t}{\tau}}$.
\end{beispiel}

\subsection{Berechnung der Faltung}

% Grafik: x(t) Rechteckimpuls [0, t₂]; h(t) Rechteckimpuls [0, t₁] und [t₁, t₂]; gespiegeltes/verschobenes x(ξ); Ausgangssignal y(t) als Trapez mit Eckpunkten bei 0, t₁, t₂

\subsection{BIBO-Stabilitätskriterium für LTI-Systeme}

$$
y(t) = \int_{-\infty}^{\infty} h(\xi)\,x(t-\xi)\,\mathrm{d}\xi
$$

Das System ist BIBO-stabil, wenn für $|x(t)| < x_{\max}$ auch $|y(t)| < \infty$ gilt:

\begin{align*}
|y(t)| &= \left|\int_{-\infty}^{\infty} h(\xi)\,x(t-\xi)\,\mathrm{d}\xi\right| \leq \int_{-\infty}^{\infty} |h(\xi)\,x(t-\xi)|\,\mathrm{d}\xi \\
       &= \int_{-\infty}^{\infty} |h(\xi)|\,|x(t-\xi)|\,\mathrm{d}\xi \leq \int_{-\infty}^{\infty} |h(\xi)|\,x_{\max}\,\mathrm{d}\xi \\
       &= x_{\max}\int_{-\infty}^{\infty} |h(\xi)|\,\mathrm{d}\xi < \infty
\end{align*}

Das System ist BIBO-stabil, wenn $\int_{-\infty}^{\infty} |h(\xi)|\,\mathrm{d}\xi < \infty$ gilt (hinreichende und notwendige Bedingung für BIBO-Stabilität).

\subsection{Die zeitdiskrete Faltung}

In gleicher Weise erhält man für ein LSI-System

\begin{align*}
y[k] = \mathcal{S}\left[x[k]\right] &= \mathcal{S}\left[\sum_{l=-\infty}^{\infty} x[l]\,\delta[k-l]\right] \\
&= \sum_{l=-\infty}^{\infty} x[l]\,\mathcal{S}\left[\delta[k-l]\right] \\
&= \sum_{l=-\infty}^{\infty} x[l]\,h[k-l] = x[k] * h[k]\,,
\end{align*}

die zeitdiskrete Faltung, so dass auch hier das Ausgangssignal $y[k]$ allein aus dem Eingangssignal $x[k]$ und der Impulsantwort $h[k]$ berechnet wird.

\subsection{Eigenschaften von LTI/LSI Systemen}

\textbf{Symmetrie} von Eingangssignal und Impulsantwort:

\begin{center}
\begin{tikzpicture}[>=latex, node distance=1cm]
  \node (x1) {$x[k]$};
  \node [draw, rectangle, right=0.8cm of x1] (h1) {$h[k]$};
  \node [right=0.8cm of h1] (y1) {$y[k]$};
  \draw[->] (x1) -- (h1); \draw[->] (h1) -- (y1);

  \node [right=1.5cm of y1] (eq) {$\equiv$};

  \node [right=1cm of eq] (x2) {$h[k]$};
  \node [draw, rectangle, right=0.8cm of x2] (h2) {$x[k]$};
  \node [right=0.8cm of h2] (y2) {$y[k]$};
  \draw[->] (x2) -- (h2); \draw[->] (h2) -- (y2);
\end{tikzpicture}
\end{center}

\textbf{Kaskadierte} Systeme sind vertauschbar:

\begin{center}
\begin{tikzpicture}[>=latex, node distance=0.8cm]
  \node (x1) {$x[k]$};
  \node [draw, rectangle, right=0.6cm of x1] (h11) {$h_1[k]$};
  \node [draw, rectangle, right=0.6cm of h11] (h12) {$h_2[k]$};
  \node [right=0.6cm of h12] (y1) {$y[k]$};
  \draw[->] (x1) -- (h11); \draw[->] (h11) -- (h12); \draw[->] (h12) -- (y1);

  \node [right=1cm of y1] (eq) {$\equiv$};

  \node [right=0.8cm of eq] (x2) {$x[k]$};
  \node [draw, rectangle, right=0.6cm of x2] (h21) {$h_2[k]$};
  \node [draw, rectangle, right=0.6cm of h21] (h22) {$h_1[k]$};
  \node [right=0.6cm of h22] (y2) {$y[k]$};
  \draw[->] (x2) -- (h21); \draw[->] (h21) -- (h22); \draw[->] (h22) -- (y2);
\end{tikzpicture}
\end{center}

\textbf{Parallele} Systeme lassen sich zusammenfassen:

\begin{center}
\begin{tikzpicture}[>=latex, node distance=0.8cm]
  \node (x) {$x[k]$};
  \node [draw, circle, right=0.6cm of x, inner sep=1pt] (sp) {$\bullet$};
  \node [draw, rectangle, above right=0.4cm and 0.6cm of sp] (h1) {$h_1[k]$};
  \node [draw, rectangle, below right=0.4cm and 0.6cm of sp] (h2) {$h_2[k]$};
  \node [draw, circle, right=2cm of sp] (sum) {$+$};
  \node [right=0.6cm of sum] (y) {$y[k]$};
  \draw[->] (x) -- (sp);
  \draw[->] (sp.north east) |- (h1); \draw[->] (h1) -| (sum);
  \draw[->] (sp.south east) |- (h2); \draw[->] (h2) -| (sum);
  \draw[->] (sum) -- (y);

  \node [right=1cm of y] (eq) {$\equiv$};

  \node [right=0.8cm of eq] (x2) {$x[k]$};
  \node [draw, rectangle, right=0.6cm of x2] (hsum) {$h_1[k]+h_2[k]$};
  \node [right=0.6cm of hsum] (y2) {$y[k]$};
  \draw[->] (x2) -- (hsum); \draw[->] (hsum) -- (y2);
\end{tikzpicture}
\end{center}

\begin{beispiel}[Schallausbreitung in halligen Räumen]
% Grafik: Raum mit Sprecher, Mikrofon y[k]; Direktschall (Freifeld), Reflexionen 1. und 2. Ordnung

\subsubsection*{Zeitdiskretes Modell der Schallausbreitung}

In einem halligen Raum gelangt der Schall nicht nur auf direktem Weg sondern auch über \textbf{Rückwürfe von den Wänden} zum Mikrofon. Auch diese verzögert eintreffenden Schallsignale werden linear überlagert. Für den Direktschall und eine Reflexion erhält man:

$$
y[k] = a_0\,x[k-k_0] + a_1\,x[k-k_1] = \sum_{l=0}^{\infty} h[l]\,x[k-l]
$$

mit

$$
h[l] = \begin{cases} a_0 & l = k_0 \\ a_1 & l = k_1 \\ 0 & \text{sonst} \end{cases}
$$

% Grafik: h[l] als zwei Impulse bei k₀ und k₁

Bei vielen und mehrfachen Reflexionen:

$$
y[k] = a_0\,x[k-k_0] + a_1\,x[k-k_1] + a_2\,x[k-k_2] + a_3\,x[k-k_3] + \ldots = \sum_{l=0}^{\infty} h[l]\,x[k-l]
$$

mit

$$
h[l] = \begin{cases} a_0 & l = k_0 \\ a_1 & l = k_1 \\ a_2 & l = k_2 \\ a_3 & l = k_3 \\ \vdots \end{cases}
$$

wobei die Überlagerung der verschiedenen Ausbreitungswege auch zu einer auslöschenden Dämpfung und damit zu negativen Koeffizienten $a_i$ führen kann.

% Grafik: h[l] mit mehreren Impulsen unterschiedlicher Vorzeichen

\subsubsection*{Raumimpulsantwort und diskrete Faltung}

\begin{itemize}
  \item Die Berücksichtigung sämtlicher Rückwürfe führt dann zur \textbf{Raumimpulsantwort} $h[l]$ mit vielen hundert von Null verschiedenen Koeffizienten.
  \item Die Impulsantwort $h[l]$ kann mit einer Impulsanregung des Raums gemessen werden.
\end{itemize}

% Grafik: Raumimpulsantwort h[l], Abtastindex l von 0 bis 2500, Amplitude −0.5 bis 0.5
\end{beispiel}

\section{Parametrische Systeme}

\subsection{Signalflussdiagramme}

\begin{center}
\begin{tabular}{lcc}
  & \textbf{Kontinuierlich} & \textbf{Diskret} \\[6pt]
  Addition: &
  \begin{tikzpicture}[>=latex, baseline=-0.5ex]
    \node (x1) {$x_1(t)$};
    \node [draw, circle, right=0.6cm of x1] (s) {$+$};
    \node [right=0.6cm of s] (y) {$y(t)$};
    \node [below=0.5cm of s] (x2) {$x_2(t)$};
    \draw[->] (x1) -- (s); \draw[->] (s) -- (y); \draw[->] (x2) -- (s);
  \end{tikzpicture}
  &
  \begin{tikzpicture}[>=latex, baseline=-0.5ex]
    \node (x1) {$x_1[k]$};
    \node [draw, circle, right=0.6cm of x1] (s) {$+$};
    \node [right=0.6cm of s] (y) {$y[k]$};
    \node [below=0.5cm of s] (x2) {$x_2[k]$};
    \draw[->] (x1) -- (s); \draw[->] (s) -- (y); \draw[->] (x2) -- (s);
  \end{tikzpicture}
  \\[18pt]
  \parbox{3cm}{Multiplikation mit\\einer Konstanten:} &
  \begin{tikzpicture}[>=latex, baseline=-0.5ex]
    \node (x) {$x(t)$};
    \node [draw, circle, right=0.6cm of x] (s) {$\otimes$};
    \node [right=0.6cm of s] (y) {$y(t)$};
    \node [below=0.5cm of s] (a) {$a$};
    \draw[->] (x) -- (s); \draw[->] (s) -- (y); \draw[->] (a) -- (s);
  \end{tikzpicture}
  &
  \begin{tikzpicture}[>=latex, baseline=-0.5ex]
    \node (x) {$x[k]$};
    \node [draw, circle, right=0.6cm of x] (s) {$\otimes$};
    \node [right=0.6cm of s] (y) {$y[k]$};
    \node [below=0.5cm of s] (a) {$a$};
    \draw[->] (x) -- (s); \draw[->] (s) -- (y); \draw[->] (a) -- (s);
  \end{tikzpicture}
  \\[18pt]
  \parbox{3cm}{Multiplikation mit\\einem Signal:} &
  \begin{tikzpicture}[>=latex, baseline=-0.5ex]
    \node (x1) {$x_1(t)$};
    \node [draw, circle, right=0.6cm of x1] (s) {$\otimes$};
    \node [right=0.6cm of s] (y) {$y(t)$};
    \node [below=0.5cm of s] (x2) {$x_2(t)$};
    \draw[->] (x1) -- (s); \draw[->] (s) -- (y); \draw[->] (x2) -- (s);
  \end{tikzpicture}
  &
  \begin{tikzpicture}[>=latex, baseline=-0.5ex]
    \node (x1) {$x_1[k]$};
    \node [draw, circle, right=0.6cm of x1] (s) {$\otimes$};
    \node [right=0.6cm of s] (y) {$y[k]$};
    \node [below=0.5cm of s] (x2) {$x_2[k]$};
    \draw[->] (x1) -- (s); \draw[->] (s) -- (y); \draw[->] (x2) -- (s);
  \end{tikzpicture}
  \\[18pt]
  Verzögerung: &
  \begin{tikzpicture}[>=latex, baseline=-0.5ex]
    \node (x) {$x(t)$};
    \node [draw, rectangle, right=0.6cm of x] (T) {$T$};
    \node [right=0.6cm of T] (y) {$y(t)$};
    \draw[->] (x) -- (T); \draw[->] (T) -- (y);
  \end{tikzpicture}
  &
  \begin{tikzpicture}[>=latex, baseline=-0.5ex]
    \node (x) {$x[k]$};
    \node [draw, rectangle, right=0.6cm of x] (T) {$T_A$};
    \node [right=0.6cm of T] (y) {$y[k]$};
    \draw[->] (x) -- (T); \draw[->] (T) -- (y);
  \end{tikzpicture}
  \\[18pt]
  Funktion: &
  \begin{tikzpicture}[>=latex, baseline=-0.5ex]
    \node (x) {$x(t)$};
    \node [draw, rectangle, right=0.6cm of x] (f) {$f(x)$};
    \node [right=0.6cm of f] (y) {$y(t)$};
    \draw[->] (x) -- (f); \draw[->] (f) -- (y);
  \end{tikzpicture}
  &
  \begin{tikzpicture}[>=latex, baseline=-0.5ex]
    \node (x) {$x[k]$};
    \node [draw, rectangle, right=0.6cm of x] (f) {$f(x)$};
    \node [right=0.6cm of f] (y) {$y[k]$};
    \draw[->] (x) -- (f); \draw[->] (f) -- (y);
  \end{tikzpicture}
\end{tabular}
\end{center}

\subsection{Parametrische LSI Systeme}

% Grafik: Vollständiges parametrisches LSI-Blockdiagramm: Vorwärtspfad x[k]→T_A→...→T_A gewichtet b₀,...,b_N; Rückkopplungspfad y[k]→T_A→...→T_A gewichtet a₁,...,a_M; alle in einer Summe

\subsection{Nichtrekursives Parametrisches LSI System (FIR)}

% Grafik: x[k] → T_A → T_A → ... → T_A, Abgriffe b₀,b₁,b₂,...,b_N, Summe → y[k]

Dieses System hat eine endliche ausgedehnte Impulsantwort und wird daher auch als ,,finite impulse-response (FIR)'' System bezeichnet.

\subsection{Rekursives Parametrisches LSI System (IIR)}

% Grafik: x[k] → b₀⊗ → (+) → y[k]; Rückkopplungspfad y[k] → T_A → T_A → ... → T_A, gewichtet a₁,a₂,a₃,...,a_M, zurück zur Summe

Dieses System hat im Allgemeinen eine unendliche ausgedehnte Impulsantwort und wird als ,,infinite impulse-response (IIR)'' System bezeichnet.

\section{Weitere Beispiele}

\begin{beispiel}[2-D Faltung in Neuronalen Netzen (CNN)]
Das Bild zeigt die grundlegende Funktion eines neuronalen Faltungsnetzes für zweidimensionale (Bild-)Daten, die durch Kreise auf der linken Seite symbolisiert sind. Die Größe des 2-D Faltungskerns (hier $M = N = 3$) ist durch das gestrichelte Quadrat angedeutet. Zur Erzeugung der Ausgangsdaten (indiziert durch $u$ und $v$) wird dann durch Verschieben des Faltungskerns eine 2-D Faltung implementiert

$$
y[u,v] = f\!\left(\sum_{m=0}^{M-1}\sum_{n=0}^{N-1} w[m,n]\,x[u-m,\,v-n] + b\right)\,,
$$

und durch einen Bias $b$ und eine nichtlineare Aktivierungsfunktion $f$ ergänzt.

% Grafik: 2-D CNN – Eingabegitter x(m,n), gestricheltes 3×3-Faltungskernelement w(m,n), Summierglied (+), Aktivierungsfunktion f → Ausgabe y(u,v)
\end{beispiel}

\begin{beispiel}[Signalübertragung mit PCM und DPCM]
Für jeden Abtastwert wird, entsprechend seiner Amplituden und der daraus resultierenden Quantisierungsstufe, ein digitaler Code übertragen. Dieses Verfahren wird deshalb \textit{pulse code modulation} (PCM) genannt.

Um die Abhängigkeit aufeinander folgender Abtastwerte für die effiziente Quantisierung zu nutzen, werden häufig differentielle Verfahren eingesetzt. Im einfachsten Fall wird vom Sender die Differenz $d(kT_A) = x(kT_A) - x((k-1)T_A)$ quantisiert, codiert und übertragen. Der Empfänger muss dann durch eine laufende Summation die Abtastwerte des Signals $x(kT_A)$ aus den Differenzwerten rekonstruieren. Mit diesem Verfahren, das als \textit{differential pulse code modulation} (DPCM) bezeichnet wird, kann Abtastrate gegen Wortbreite abgetauscht werden.

\subsubsection*{DPCM Sender}

\begin{center}
\begin{tikzpicture}[>=latex, node distance=1cm]
  \node (xt) {$x(t)$};
  \node [draw, rectangle, right=0.8cm of xt] (abt) {Abtastung};
  \node [draw, circle, right=0.8cm of abt] (sum) {$+$};
  \node [right=0.8cm of sum] (dk) {$d[k]$};
  \node [draw, rectangle, below=1cm of sum, text width=1.4cm, align=center] (zsp) {Zwischen\-speicher};
  \node [below=0.5cm of xt] (TA) {$T_A$};

  \draw[->] (xt) -- (abt);
  \draw[->] (abt) -- node[above] {$x[k]$} (sum);
  \draw[->] (sum) -- (dk);
  \draw[->] (TA) -- (abt);
  \draw[->] (sum.south) -- (zsp.north);
  \draw[->] (zsp.west) -| node[left, near end] {$z_s[k]$} (sum.south west);
\end{tikzpicture}
\end{center}

Berechnung der zeitdiskreten Ausgangssignale für $k \geq 0$ mit Differenzengleichungen. Für den Sender erhält man:

$$
d[k] = x[k] - z_s[k] = \begin{cases} x[0] - z_s[0] & k = 0 \\ x[k] - x[k-1] & k \geq 1 \end{cases}
$$

\subsubsection*{DPCM Empfänger}

\begin{center}
\begin{tikzpicture}[>=latex, node distance=1cm]
  \node (dk) {$d[k]$};
  \node [draw, circle, right=0.8cm of dk] (sum) {$+$};
  \node [right=0.8cm of sum] (xhk) {$\widehat{x}[k]$};
  \node [draw, rectangle, right=1cm of xhk] (rek) {Rekonstruktion};
  \node [right=0.8cm of rek] (xht) {$\widehat{x}(t)$};
  \node [draw, rectangle, below=1cm of sum, text width=1.4cm, align=center] (zsp) {Zwischen\-speicher};
  \node [below=0.5cm of rek] (TA) {$T_A$};

  \draw[->] (dk) -- (sum);
  \draw[->] (sum) -- (xhk);
  \draw[->] (xhk) -- (rek);
  \draw[->] (rek) -- (xht);
  \draw[->] (TA) -- (rek);
  \draw[->] (xhk.south) |- (zsp.east);
  \draw[->] (zsp.west) -| node[left, near end] {$z_E[k]$} (sum.south west);
\end{tikzpicture}
\end{center}

Für den Empfänger erhält man:

$$
\widehat{x}[k] = d[k] + z_E[k] = z_E[0] + \sum_{n=0}^{k} d[n]
$$

Für $z_S[0] = z_E[0]$ gilt $\widehat{x}[k] = x[k]\ \forall\ k \geq 0$. Sender und Empfänger enthalten Speicher für einen Abtastwert bzw.\ die laufende Summe. Diese stellen die inneren Zustände dar.

\subsubsection*{Differenzsignale}

% Grafik: Zwei Plots bei f_A = 16 kHz; oben x(kT_A) (Originalsignal), unten d(kT_A) (Differenzsignal), t/s von 0 bis 6

% Grafik: Zwei Plots bei f_A = 64 kHz; oben x(kT_A), unten d(kT_A), t/s von 0 bis 6

\subsubsection*{Original und Differenzsignal}

% Grafik (4 Plots): Vergleich f_A = 16 kHz und f_A = 64 kHz, jeweils x(kT_A) und d(kT_A) im Ausschnitt kT_A von 0.156562 bis ~0.1578

\subsubsection*{Differenzsignal und rekonstruiertes Signal}

% Grafik (4 Plots): f_A = 16 kHz und f_A = 64 kHz – oben d(kT_A), unten x̂(kT_A), Ausschnitt kT_A ≈ 0.156562–0.1578

\subsubsection*{Quantisiertes (2 Bit) und rekonstruiertes Signal}

% Grafik (6 Plots): f_A = 16 kHz und f_A = 64 kHz – oben Q_q(d(kT_A)) (2-Bit-quantisiertes Differenzsignal), Mitte x̂(kT_A) (rekonstruiert), unten e²(kT_A) (quadratischer Fehler, Skala ×10⁻³)
\end{beispiel}

\section*{Zusammenfassung}

\begin{itemize}
  \item Systeme bilden Eingangssignale auf Ausgangssignale ab.
  \item Hierbei spielen die inneren Zustandvariablen des Systems eine große Rolle.
  \item Wichtige Systemeigenschaften sind die Linearität, die Zeitinvarianz oder Verschiebungsinvarianz und die Stabilität.
  \item LTI bzw.\ LSI Systeme werden durch ihre Impulsantwort beschrieben.
  \item Zeitinvariante und zeitvariante Systeme können mit einem Zustandsraummodell beschrieben und berechnet werden.
\end{itemize}


